Contextual retrieval prepends a short document-level summary to every chunk
before indexing, and the published recipe pairs it with a cross-encoder reranker
to restore the precision that flat chunks lose. Neither stage is cheap outside a
datacenter, and this paper asks whether an edge pipeline needs both. We label
3{,}134 chunks from 82 Russian business documents with an API teacher (24.3M
prompt tokens, USD~0.82) and distill the teacher's contextual enrichment into
\texttt{Qwen3-0.6B} with LoRA. Retrieval is then measured with
\texttt{Qwen3-Embedding-0.6B}, BM25, reciprocal rank fusion and
\texttt{Qwen3-Reranker-0.6B} on an Apple~M2 GPU (MLX, fp16): six configurations
on 234 synthetic queries that carry gold chunks, and all eight on a stratified
subset of sixty. On that subset a reranker-free pipeline with distilled
on-device context reaches 0.600 recall@10 against 0.650 for the same index with
a 0.6B cross-encoder, a difference whose paired bootstrap interval includes
zero, at a median latency of 41.4~ms against 54.7~s, a factor of 1{,}322, while
reranking flat chunks is dominated by the contextual dense-only arm. On this
homogeneous corpus BM25 fusion degrades quality and no setting of the fusion
parameters recovers the loss. Distilling the context generator costs little in
quality: the student lands within 0.04 recall@10 of the teacher while improving
recall@1 and MRR, and with Matryoshka truncation to 512 dimensions plus a
Russian query instruction it matches the teacher baseline on recall@1 and MRR at
half the index size. Generating the contexts for the whole corpus takes 1.85~h
and USD~0.82 through the API, or 4.49~h at no cost on the device, where no
document leaves the machine. The dataset, the adapter, the evaluation harness
and the measurement scripts are released (Section~\ref{sec:availability}).
